% ID: SEL-D0008-BC-E06
% ESTADO: cotejada. Segunda pasada visual completa y borrador compilado revisados.
% Fuente: Selectividad/matematicas_aplicadas/ORD_SUPLENTE_Modelo-4-Enunciados (Mat-Ap-2021).pdf
% SHA-256: 74b067b81be20aa0262a0cc58795f3b25488aa9c749d825b7b33f656dd1ae089
% Página: 2; bloque C; ejercicio original 6.

Una urna $A$ contiene 4 bolas rojas y 5 verdes y otra urna $B$ contiene
6 bolas rojas y 3 verdes. Lanzamos dos dados y si la suma es mayor o igual
a 9, extraemos una bola de la urna $A$ y en caso contrario, la extraemos
de la urna $B$.
\begin{enumerate}[label=\alph*)]
  \item \textbf{(1.5 puntos)} Calcule la probabilidad de que la bola extraída
  sea verde y de la urna $B$.
  \item \textbf{(1 punto)} Halle la probabilidad de que la bola extraída
  sea roja.
\end{enumerate}
